\olchapter{fol}{tab}{Tableaux}

\begin{editorial}
This chapter presents a signed analytic tableaux system.
  
To include or exclude material relevant to natural deduction as a
proof system, use the ``prfTab'' tag.
\end{editorial}



\olfileid{fol}{tab}{rul}

\olsection{Rules and \usetoken{P}{tableau}}

!!^a{tableau} is a systematic survey of the possible ways
!!a{sentence} can be true or false in !!a{structure}. The building
blocks of a tableau are !!{signed formula}s: !!{sentence}s plus a
truth value ``sign,'' either $\True$ or~$\False$. These signed
!!{formula}s are arranged in a (downward growing) tree.

\begin{defn}
  A \emph{!!{signed formula}} is a pair consisting of a truth value
  and !!a{sentence}, i.e., either:
  \[
  \sFmla{\True}{!A} \text{ or } \sFmla{\False}{!A}.
  \]
\end{defn}

Intuitively, we might read $\sFmla{\True}{!A}$ as ``$!A$ might be
true'' and $\sFmla{\False}{!A}$ as ``$!A$ might be false'' (in some
!!{structure}).

Each !!{signed formula} in the tree is either an \emph{assumption}
(which are listed at the very top of the tree), or it is obtained from
!!a{signed formula} above it by one of a number of rules of
inference. There are two rules for each possible !!{main operator} of
the preceding !!{formula}, one for the case where the sign is~$\True$,
and one for the case where the sign is~$\False$. Some rules allow the
tree to branch, and some only add !!{signed formula}s to the branch.
A rule may be (and often must be) applied not to the immediately
preceding !!{signed formula}, but to any !!{signed formula} in the
branch from the root to the place the rule is applied.

A branch is \emph{closed} when it contains both $\sFmla{\True}{!A}$
and $\sFmla{\False}{!A}$. A closed !!{tableau} is one where every branch
is closed.  Under the intuitive interpretation, any branch describes a
joint possibility, but $\sFmla{\True}{!A}$ and $\sFmla{\False}{!A}$
are not jointly possible. In other words, if a branch is closed, the
possibility it describes has been ruled out. In particular, that means
that a closed !!{tableau} rules out all possibilities of simultaneously
making every assumption of the form $\sFmla{\True}{!A}$ true and every
assumption of the form~$\sFmla{\False}{!A}$ false.

A closed !!{tableau} \emph{for $!A$} is a closed !!{tableau} with
root~$\sFmla{\False}{!A}$. If such a closed !!{tableau} exists, all
possibilities for~$!A$ being false have been ruled out; i.e., $!A$
must be true in every !!{structure}.





\olfileid{fol}{tab}{prl}

\olsection{Propositional Rules}

\subsection{Rules for $\lnot$}

\begin{defish}
\AxiomC{\sFmla{\True}{\lnot !A}}
\RightLabel{\TRule{\True}{\lnot}}
\UnaryInfC{\sFmla{\False}{!A}}
\DisplayProof
\hfill
\AxiomC{\sFmla{\False}{\lnot !A}}
\RightLabel{\TRule{\False}{\lnot}}
\UnaryInfC{\sFmla{\True}{!A}}
\DisplayProof
\end{defish}

\subsection{Rules for $\land$}

\begin{defish}\noindent
\AxiomC{\sFmla{\True}{!A \land !B}}
\RightLabel{\TRule{\True}{\land}}
\UnaryInfC{\sFmla{\True}{!A}}
\noLine
\UnaryInfC{\sFmla{\True}{!B}}
\DisplayProof
\hfill
\AxiomC{\sFmla{\False}{!A \land !B}}
\RightLabel{\TRule{\False}{\land}}
\UnaryInfC{$\sFmla{\False}{!A} \quad \mid \quad \sFmla{\False}{!B}$}
\DisplayProof
\end{defish}

\subsection{Rules for $\lor$}

\begin{defish}
\AxiomC{\sFmla{\True}{!A \lor !B}}
\RightLabel{\TRule{\True}{\lor}}
\UnaryInfC{$\sFmla{\True}{!A} \quad \mid \quad \sFmla{\True}{!B}$}
\DisplayProof
\hfill
\AxiomC{\sFmla{\False}{!A \lor !B}}
\RightLabel{\TRule{\False}{\lor}}
\UnaryInfC{\sFmla{\False}{!A}}
\noLine
\UnaryInfC{\sFmla{\False}{!B}}
\DisplayProof
\end{defish}

\subsection{Rules for $\lif$}

\begin{defish}
\AxiomC{\sFmla{\True}{!A \lif !B}}
\RightLabel{\TRule{\True}{\lif}}
\UnaryInfC{$\sFmla{\False}{!A} \quad \mid \quad \sFmla{\True}{!B}$}
\DisplayProof
\hfill
\AxiomC{\sFmla{\False}{!A \lif !B}}
\RightLabel{\TRule{\False}{\lif}}
\UnaryInfC{\sFmla{\True}{!A}}
\noLine
\UnaryInfC{\sFmla{\False}{!B}}
\DisplayProof
\end{defish}

\subsection{The Cut Rule}

\begin{defish}
\AxiomC{}
\RightLabel{\Cut}
\UnaryInfC{$\sFmla{\True}{!A} \quad \mid \quad \sFmla{\False}{!A}$}
\DisplayProof
\end{defish}

The \Cut{} rule is not applied ``to'' a previous !!{signed formula}; rather,
it allows every branch in !!a{tableau} to be split in two, one branch
containing $\sFmla{\True}{!A}$, the other~$\sFmla{\False}{!A}$. It is
not necessary---any set of !!{signed formula}s with a closed
!!{tableau} has one not using \Cut---but it allows us to combine
!!{tableau}s in a convenient way.




  

\olfileid{fol}{tab}{qrl}

\olsection{Quantifier Rules}

\subsection{Rules for $\lforall$}

\begin{defish}
\AxiomC{\sFmla{\True}{\lforall[x][!A(x)]}}
\RightLabel{\TRule{\True}{\forall}}
\UnaryInfC{\sFmla{\True}{!A(t)}}
\DisplayProof
\hfill
\AxiomC{\sFmla{\False}{\lforall[x][!A(x)]}}
\RightLabel{\TRule{\False}{\lforall}}
\UnaryInfC{\sFmla{\False}{!A(a)}}
\DisplayProof
\end{defish}

In \TRule{\True}{\lforall}, $t$ is a closed term (i.e., one without
variables). In \TRule{\False}{\lforall}, $a$~is !!a{constant} which
must not occur anywhere in the branch above \TRule{\False}{\lforall}
rule. We call $a$ the \emph{eigenvariable} of the
\TRule{\False}{\forall} inference.\footnote{We use the term
``eigenvariable'' even though $a$ in the above rule is !!a{constant}.
This has historical reasons.}

\subsection{Rules for $\lexists$}

\begin{defish}
\AxiomC{\sFmla{\True}{\lexists[x][!A(x)]}}
\RightLabel{\TRule{\True}{\lexists}}
\UnaryInfC{\sFmla{\True}{!A(a)}}
\DisplayProof
\hfill
\AxiomC{\sFmla{\False}{\lexists[x][!A(x)]}}
\RightLabel{\TRule{\False}{\lexists}}
\UnaryInfC{\sFmla{\False}{!A(t)}}
\DisplayProof
\end{defish}

Again, $t$~is a closed term, and $a$~is !!a{constant} which does not
occur in the branch above the~\TRule{\True}{\lexists} rule. We call
$a$ the \emph{eigenvariable} of the \TRule{\True}{\lexists} inference.

The condition that an eigenvariable not occur in the branch above the
\TRule{\False}{\lforall} or \TRule{\True}{\lexists} inference is
called the \emph{eigenvariable condition}.

\begin{explain}
Recall the convention that when $!A$ is !!a{formula} with the
!!{variable}~$x$ free, we indicate this by writing~$!A(x)$. In the
same context, $!A(t)$ then is short for~$\Subst{!A}{t}{x}$. So we
could also write the \TRule{\False}{\lexists} rule as:
\begin{prooftree}
  \AxiomC{\sFmla{\False}{\lexists[x][!A]}}
  \RightLabel{\TRule{\False}{\lexists}}
  \UnaryInfC{\sFmla{\False}{\Subst{!A}{t}{x}}}
\end{prooftree}
Note that $t$ may already occur in~$!A$, e.g., $!A$~might
be~$\Atom{\Obj P}{t,x}$. Thus, inferring $\sFmla{\False}{\Atom{\Obj
P}{t,t}}$ from~$ \sFmla{\False}{\lexists[x][\Atom{\Obj P}{t,x}]}$ is
a correct application of~\TRule{\False}{\lexists}. However, the
eigenvariable conditions in \TRule{\False}{\lforall}
and~\TRule{\True}{\lexists} require that the !!{constant}~$a$ does
not occur in~$!A$. So, you cannot correctly infer
$\sFmla{\False}{\Atom{\Obj P}{a,a}}$ from
$\sFmla{\False}{\lforall[x][\Atom{\Obj P}{a,x}]}$
using~$\TRule{\False}{\lforall}$.
\end{explain}

\begin{explain}
In \TRule{\True}{\lforall} and \TRule{\False}{\lexists} there are no
restrictions on the term~$t$. On the other hand, in the
\TRule{\True}{\lexists} and \TRule{\False}{\lforall} rules, the
eigenvariable condition requires that the !!{constant}~$a$ does not
occur anywhere in the branches above the respective inference. It is
necessary to ensure that the system is sound. Without this condition,
the following would be a closed !!{tableau} for
$\lexists[x][\formula{A}(x)] \lif \lforall[x][\formula{A}(x)]$:
\begin{center}
\begin{tableau}{}
  [\sFmla{\False}{\lexists[x][\formula{A}(x)] \lif \lforall[x][\formula{A}(x)]}, just=\TAss
    [\sFmla{\True}{\lexists[x][\formula{A}(x)]},
      just={\TRule{\False}{\lif}[1]}
      [\sFmla{\False}{\lforall[x][\formula{A}(x)]},
        just={\TRule{\False}{\lif}[1]}
        [\sFmla{\True}{\formula{A}(a)},
          just={\TRule{\True}{\lexists}[2]}
          [\sFmla{\False}{\formula{A}(a)},
            just={\TRule{\False}{\lforall}[3]}, close]
        ]
      ]
    ]
  ]
\end{tableau}
\end{center}
However, $\lexists[x][\formula{A}(x)] \lif
\lforall[x][\formula{A}(x)]$ is not valid.
\end{explain}






\olfileid{fol}{tab}{der}

\olsection{\usetoken{P}{tableau}}

\begin{explain}
We've said what an assumption is, and we've given the rules of
inference.  !!^{tableau}s are inductively generated from these: each
!!{tableau} either is a single branch consisting of one or more
assumptions, or it results from !!a{tableau} by applying one of the
rules of inference on a branch.
\end{explain}

\begin{defn}[!!^{tableau}]
!!^a{tableau} for assumptions $\sFmla{S_1}{!A_1}$, \dots,
$\sFmla{S_n}{!A_n}$ (where each $S_i$ is either $\True$ or~$\False$) is
a finite tree of !!{signed formula}s satisfying the following conditions:
\begin{enumerate}
\item The $n$ topmost !!{signed formula}s of the tree are
  $\sFmla{S_i}{!A_i}$, one below the other.
\item Every !!{signed formula} in the tree that is not one of the
  assumptions results from a correct application of an inference rule
  to !!a{signed formula} in the branch above it.
\end{enumerate}
A branch of !!a{tableau} is \emph{closed} iff it contains both
$\sFmla{\True}{!A}$ and~$\sFmla{\False}{!A}$, and \emph{open}
otherwise. !!^a{tableau} in which every branch is closed is a
\emph{closed !!{tableau}} (for its set of assumptions). If !!a{tableau} is
not closed, i.e., if it contains at least one open branch, it is
\emph{open}.
\end{defn}

\begin{ex}
  Every set of assumptions on its own is !!a{tableau}, but it will
  generally not be closed. (Obviously, it is closed only if the
  assumptions already contain a pair of !!{signed formula}s
  $\sFmla{\True}{!A}$ and~$\sFmla{\False}{!A}$.)

  From !!a{tableau} (open or closed) we can obtain a new, larger one by
  applying one of the rules of inference to !!a{signed formula}~$!A$
  in it. The rule will append one or more !!{signed formula}s to the end of
  any branch containing the occurrence of~$!A$ to which we apply the
  rule.

  For instance, consider the assumption $\sFmla{\True}{!A \land \lnot
    !A}$. Here is the (open) !!{tableau} consisting of just that
  assumption:
  \begin{center}
    \begin{tableau}{}
      [\sFmla{\True}{\formula{A} \land \lnot \formula{A}}, just=\TAss]
    \end{tableau}{}
  \end{center}
  We obtain a new !!{tableau} from it by applying the $\TRule{\True}{\land}$
  rule to the assumption. That rule allows us to add two new lines to
  the !!{tableau}, $\sFmla{\True}{!A}$ and $\sFmla{\True}{\lnot !A}$:
  \begin{center}
    \begin{tableau}{}
      [\sFmla{\True}{\formula{A} \land \lnot \formula{A}}, just=\TAss
        [\sFmla{\True}{\formula{A}}, just={\TRule{\True}{\land}[1]},
          [\sFmla{\True}{\lnot \formula{A}}, just={\TRule{\True}{\land}[1]}
          ]
        ]
      ]
    \end{tableau}{}
  \end{center}
  When we write down !!{tableau}s, we record the rules we've applied
  on the right (e.g., $\TRule{\True}{\land} 1$ means that the
  !!{signed formula} on that line is the result of applying the
  $\TRule{\True}{\land}$ rule to the !!{signed formula} on line~$1$).
  This new !!{tableau} now contains additional !!{signed formula}s,
  but to only one ($\sFmla{\True}{\lnot !A}$) can we apply a rule (in
  this case, the $\TRule{\True}{\lnot}$ rule). This results in the closed
  !!{tableau}
  \begin{center}
    \begin{tableau}{}
      [\sFmla{\True}{\formula{A} \land \lnot \formula{A}}, just=\TAss
        [\sFmla{\True}{\formula{A}}, just={\TRule{\True}{\land}[1]}
          [\sFmla{\True}{\lnot \formula{A}}, just={\TRule{\True}{\land}[1]}
            [\sFmla{\False}{\formula{A}}, just={\TRule{\True}{\lnot}[3]}, close]
          ]
        ]
      ]
    \end{tableau}{}
  \end{center}
\end{ex}





\olfileid{fol}{tab}{pro}

\olsection{Examples of \usetoken{P}{tableau}}

\begin{ex}
Let's find a closed !!{tableau} for the !!{sentence} $(!A \land !B) \lif !A$.

We begin by writing the corresponding assumption at the top of the
!!{tableau}.
\begin{oltableau}
  [\sFmla{\False}{(\formula{A} \land \formula{B}) \lif \formula{A}},
    just = \TAss]
\end{oltableau}

There is only one assumption, so only one !!{signed formula} to which
we can apply a rule. (For every !!{signed formula}, there is always at
most one rule that can be applied: it's the rule for the corresponding
sign and !!{main operator} of the !!{sentence}.) In this case, this
means, we must apply $\TRule{\False}{\lif}$.
\begin{oltableau}
  [\sFmla{\False}{(\formula{A} \land \formula{B}) \lif \formula{A}},
    checked, just = \TAss
    [\sFmla{\True}{\formula{A} \land \formula{B}},
      just={\TRule{\False}{\lif}[1]}
      [\sFmla{\False}{\formula{A}}, just={\TRule{\False}{\lif}[1]}]
    ]
  ]
\end{oltableau}
To keep track of which !!{signed formula}s we have applied their
corresponding rules to, we write a checkmark next to the
sentence. However, \emph{only} write a checkmark if the rule has been
applied to all open branches. Once !!a{signed formula} has had the
corresponding rule applied in every open branch, we will not have to
return to it and apply the rule again. In this case, there is only one
branch, so the rule only has to be applied once. (Note that checkmarks
are only a convenience for constructing tableaux and are not
officially part of the syntax of tableaux.)

There is one new !!{signed formula} to which we can apply a rule: the
$\sFmla{\True}{!A \land !B}$ on line~$2$. Applying the
$\TRule{\True}{\land}$ rule results in:
\begin{oltableau}
  [\sFmla{\False}{(\formula{A} \land \formula{B}) \lif \formula{A}},
    checked, just = \TAss
    [\sFmla{\True}{\formula{A} \land \formula{B}},
      just={\TRule{\False}{\lif}[1]}, checked
      [\sFmla{\False}{\formula{A}}, just={\TRule{\False}{\lif}[1]}
        [\sFmla{\True}{\formula{A}}, just={\TRule{\True}{\land}[2]}
          [\sFmla{\True}{\formula{B}}, just={\TRule{\True}{\land}[2]}, close
          ]
        ]
      ]
    ]
  ]
\end{oltableau}
Since the branch now contains both $\sFmla{\True}{!A}$ (on line~$4$)
and $\sFmla{\False}{!A}$ (on line~$3$), the branch is closed. Since it
is the only branch, the !!{tableau} is closed. We have found
a closed !!{tableau} for~$(!A \land !B) \lif !A$.
\end{ex}

\begin{ex}
Now let's find a closed !!{tableau} for $(\lnot !A \lor !B) \lif (!A
\lif !B)$.

We begin with the corresponding assumption:
\begin{oltableau}
  [\sFmla{\False}{(\lnot \formula{A} \lor \formula{B}) \lif
      (\formula{A} \lif \formula{B})}, just=\TAss]
\end{oltableau}
The one !!{signed formula} in this !!{tableau} has !!{main operator}~$\lif$
and sign~$\False$, so we apply the $\TRule{\False}{\lif}$ rule to it
to obtain:
\begin{oltableau}
  [\sFmla{\False}{(\lnot \formula{A} \lor \formula{B})
      \lif (\formula{A} \lif \formula{B})}, just=\TAss, checked
    [\sFmla{\True}{\lnot \formula{A} \lor \formula{B}},
      just={\TRule{\False}{\lif}[1]}
      [\sFmla{\False}{(\formula{A} \lif \formula{B})},
        just={\TRule{\False}{\lif}[1]}
      ]
    ]
  ]
\end{oltableau}
We now have a choice as to whether to apply~$\TRule{\True}{\lor}$ to
line~$2$ or $\TRule{\False}{\lif}$ to line~$3$. It actually doesn't
matter which order we pick, as long as each !!{signed formula} has its
corresponding rule applied in every branch. So let's pick the first
one. The $\TRule{\True}{\lor}$ rule allows the !!{tableau} to branch,
and the two conclusions of the rule will be the new !!{signed formula}s
added to the two new branches. This results in:
\begin{oltableau}
  [\sFmla{\False}{(\lnot \formula{A} \lor \formula{B}) \lif
      (\formula{A} \lif \formula{B})}, just=\TAss, checked
    [\sFmla{\True}{\lnot \formula{A} \lor \formula{B}},
      just={\TRule{\False}{\lif}[1]}, checked
      [\sFmla{\False}{(\formula{A} \lif \formula{B})},
        just={\TRule{\False}{\lif}[1]}
        [\sFmla{\True}{\lnot \formula{A}}, just={\TRule{\True}{\lor}[2]}]
        [\sFmla{\True}{\formula{B}}, just={\TRule{\True}{\lor}[2]}]
      ]
    ]
  ]
\end{oltableau}
We have not applied the $\TRule{\False}{\lif}$ rule to line~$3$ yet:
let's do that now.  To save time, we apply it to both branches.
Recall that we write a checkmark next to !!a{signed formula} only if
we have applied the corresponding rule in every open branch. So it's a
good idea to apply a rule at the end of every branch that contains the
!!{signed formula} the rule applies to. That way we won't have to
return to that !!{signed formula} lower down in the various branches.
\begin{oltableau}
  [\sFmla{\False}{(\lnot \formula{A} \lor \formula{B}) \lif
      (\formula{A} \lif \formula{B})}, just=\TAss, checked
    [\sFmla{\True}{\lnot \formula{A} \lor \formula{B}},
      just={\TRule{\False}{\lif}[1]}, checked
      [\sFmla{\False}{(\formula{A} \lif \formula{B})},
        just={\TRule{\False}{\lif}[1]}, checked
        [\sFmla{\True}{\lnot \formula{A}}, just={\TRule{\True}{\lor}[2]}
          [\sFmla{\True}{\formula{A}}, just={\TRule{\False}{\lif}[3]}
            [\sFmla{\False}{\formula{B}}, just={\TRule{\False}{\lif}[3]}]
          ]
        ]
        [\sFmla{\True}{\formula{B}}, just={\TRule{\True}{\lor}[2]}
          [\sFmla{\True}{\formula{A}}, just={\TRule{\False}{\lif}[3]}
            [\sFmla{\False}{\formula{B}}, just={\TRule{\False}{\lif}[3]}, close]
          ]
        ]
      ]
    ]
  ]
\end{oltableau}
The right branch is now closed. On the left branch, we can still apply
the $\TRule{\True}{\lnot}$ rule to line~$4$. This results
in~$\sFmla{\False}{!A}$ and closes the left branch:
\begin{oltableau}
  [\sFmla{\False}{(\lnot \formula{A} \lor \formula{B}) \lif
      (\formula{A} \lif \formula{B})}, just=\TAss, checked
    [\sFmla{\True}{\lnot \formula{A} \lor \formula{B}},
      just={\TRule{\False}{\lif}[1]}, checked
      [\sFmla{\False}{(\formula{A} \lif \formula{B})},
        just={\TRule{\False}{\lif}[1]}, checked
        [\sFmla{\True}{\lnot \formula{A}}, just={\TRule{\True}{\lor}[2]}
          [\sFmla{\True}{\formula{A}}, just={\TRule{\False}{\lif}[3]}
            [\sFmla{\False}{\formula{B}}, just={\TRule{\False}{\lif}[3]}
              [\sFmla{\False}{\formula{A}},
                just={\TRule{\True}{\lnot}[4]}, close]
            ]
          ]
        ]
        [\sFmla{\True}{\formula{B}}, just={\TRule{\True}{\lor}[2]}
          [\sFmla{\True}{\formula{A}}, just={\TRule{\False}{\lif}[3]}
            [\sFmla{\False}{\formula{B}},
              just={\TRule{\False}{\lif}[3]}, close]
          ]
        ]
      ]
    ]
  ]
\end{oltableau}
\end{ex}

\begin{ex}
We can give !!{tableau}s for any number of !!{signed formula}s as
assumptions.  Often it is also necessary to apply more than one rule
that allows branching; and in general !!a{tableau} can have any number
of branches. For instance, consider !!a{tableau} for
$\{\sFmla{\True}{!A \lor (!B \land !C)}, \sFmla{\False}{(!A \lor !B) \land (!A
  \lor !C)}\}$. We start by applying the $\TRule{\True}{\lor}$ to the
first assumption:
\begin{oltableau}
  [\sFmla{\True}{\formula{A} \lor (\formula{B} \land \formula{C})},
    just=\TAss, checked
    [\sFmla{\False}{(\formula{A} \lor \formula{B}) \land
        (\formula{A} \lor \formula{C})}, just=\TAss
      [\sFmla{\True}{\formula{A}}, just={\TRule{\True}{\lor}[1]}]
      [\sFmla{\True}{\formula{B} \land \formula{C}},
        just={\TRule{\True}{\lor}[1]}]
    ]
  ]
\end{oltableau}
Now we can apply the $\TRule{\False}{\land}$ rule to line~$2$. We do
this on both branches simultaneously, and can therefore check off
line~$2$:
\begin{oltableau}
  [\sFmla{\True}{\formula{A} \lor (\formula{B} \land \formula{C})},
    just=\TAss, checked
    [\sFmla{\False}{(\formula{A} \lor \formula{B}) \land
        (\formula{A} \lor \formula{C})}, just=\TAss, checked
      [\sFmla{\True}{\formula{A}}, just={\TRule{\True}{\lor}[1]}
        [\sFmla{\False}{\formula{A} \lor \formula{B}},
          just={\TRule{\False}{\land}[2]}]
        [\sFmla{\False}{\formula{A} \lor \formula{C}},
          just={\TRule{\False}{\land}[2]}]
      ]
      [\sFmla{\True}{\formula{B} \land \formula{C}},
        just={\TRule{\True}{\lor}[1]}
        [\sFmla{\False}{\formula{A} \lor \formula{B}},
          just={\TRule{\False}{\land}[2]}]
        [\sFmla{\False}{\formula{A} \lor \formula{C}},
          just={\TRule{\False}{\land}[2]}]
      ]
    ]
  ]
\end{oltableau}
Now we can apply $\TRule{\False}{\lor}$ to all the branches containing
$!A \lor !B$:
\begin{oltableau}
  [\sFmla{\True}{\formula{A} \lor (\formula{B} \land \formula{C})},
    just=\TAss, checked
    [\sFmla{\False}{(\formula{A} \lor \formula{B}) \land
        (\formula{A} \lor \formula{C})}, just=\TAss, checked
      [\sFmla{\True}{\formula{A}}, just={\TRule{\True}{\lor}[1]}
        [\sFmla{\False}{\formula{A} \lor \formula{B}},
          just={\TRule{\False}{\land}[2]}, checked
          [\sFmla{\False}{\formula{A}}, just={\TRule{\False}{\lor}[4]}
            [\sFmla{\False}{\formula{B}},
              just={\TRule{\False}{\lor}[4]}, close]
          ]
        ]
        [\sFmla{\False}{\formula{A} \lor \formula{C}},
          just={\TRule{\False}{\land}[2]}]
      ]
      [\sFmla{\True}{\formula{B} \land \formula{C}},
        just={\TRule{\True}{\lor}[1]}
        [\sFmla{\False}{\formula{A} \lor \formula{B}},
          just={\TRule{\False}{\land}[2]}, checked
          [\sFmla{\False}{\formula{A}}, just={\TRule{\False}{\lor}[4]}
            [\sFmla{\False}{\formula{B}},
              just={\TRule{\False}{\lor}[4]}]
          ]
        ]
        [\sFmla{\False}{\formula{A} \lor \formula{C}},
          just={\TRule{\False}{\land}[2]}]
      ]
    ]
  ]
\end{oltableau}
The leftmost branch is now closed. Let's now apply
$\TRule{\False}{\lor}$ to $!A \lor !C$:
\begin{oltableau}
  [\sFmla{\True}{\formula{A} \lor (\formula{B} \land \formula{C})},
    just=\TAss, checked
    [\sFmla{\False}{(\formula{A} \lor \formula{B}) \land
        (\formula{A} \lor \formula{C})}, just=\TAss, checked
      [\sFmla{\True}{\formula{A}}, just={\TRule{\True}{\lor}[1]}
        [\sFmla{\False}{\formula{A} \lor \formula{B}},
          just={\TRule{\False}{\land}[2]}, checked
          [\sFmla{\False}{\formula{A}},
            just={\TRule{\False}{\lor}[4]}
            [\sFmla{\False}{\formula{B}},
              just={\TRule{\False}{\lor}[4]}, close]
          ]
        ]
        [\sFmla{\False}{\formula{A} \lor \formula{C}},
          just={\TRule{\False}{\land}[2]}, checked
          [\sFmla{\False}{\formula{A}},
            just={\TRule{\False}{\lor}[4]}, move by=2
            [\sFmla{\False}{\formula{C}},
              just={\TRule{\False}{\lor}[4]}, close]
          ]
        ]
      ]
      [\sFmla{\True}{\formula{B} \land \formula{C}},
        just={\TRule{\True}{\lor}[1]}
        [\sFmla{\False}{\formula{A} \lor \formula{B}},
          just={\TRule{\False}{\land}[2]}, checked
          [\sFmla{\False}{\formula{A}},
            just={\TRule{\False}{\lor}[4]}
            [\sFmla{\False}{\formula{B}},
              just={\TRule{\False}{\lor}[4]}
            ]
          ]
        ]
        [\sFmla{\False}{\formula{A} \lor \formula{C}},
          just={\TRule{\False}{\land}[2]}, checked
          [\sFmla{\False}{\formula{A}},
            just={\TRule{\False}{\lor}[4]}, move by=2
            [\sFmla{\False}{\formula{C}},
              just={\TRule{\False}{\lor}[4]}]
          ]
        ]
      ]
    ]
  ]
\end{oltableau}
Note that we moved the result of applying $\TRule{\False}{\lor}$ a
second time below for clarity. In this instance it would not have been
needed, since the justifications would have been the same.

Two branches remain open, and $\sFmla{\True}{!B \land !C}$ on line~$3$
remains unchecked. We apply $\TRule{\True}{\land}$ to it to obtain a
closed !!{tableau}:
\begin{oltableau}
  [\sFmla{\True}{\formula{A} \lor (\formula{B} \land \formula{C})},
    just=\TAss, checked
    [\sFmla{\False}{(\formula{A} \lor \formula{B}) \land
        (\formula{A} \lor \formula{C})}, just=\TAss, checked
      [\sFmla{\True}{\formula{A}}, just={\TRule{\True}{\lor}[1]}
        [\sFmla{\False}{\formula{A} \lor \formula{B}},
          just={\TRule{\False}{\land}[2]}, checked
          [\sFmla{\False}{\formula{A}},
            just={\TRule{\False}{\lor}[4]}
            [\sFmla{\False}{\formula{B}},
              just={\TRule{\False}{\lor}[4]}, close]
          ]
        ]
        [\sFmla{\False}{\formula{A} \lor \formula{C}},
          just={\TRule{\False}{\land}[2]}, checked
          [\sFmla{\False}{\formula{A}},
            just={\TRule{\False}{\lor}[4]}
            [\sFmla{\False}{\formula{C}},
              just={\TRule{\False}{\lor}[4]}, close]
          ]
        ]
      ]
      [\sFmla{\True}{\formula{B} \land \formula{C}},
        just={\TRule{\True}{\lor}[1]}, checked
        [\sFmla{\False}{\formula{A} \lor \formula{B}},
          just={\TRule{\False}{\land}[2]}, checked
          [\sFmla{\False}{\formula{A}},
            just={\TRule{\False}{\lor}[4]}
            [\sFmla{\False}{\formula{B}},
              just={\TRule{\False}{\lor}[4]}
              [\sFmla{\True}{\formula{B}},
                just={\TRule{\True}{\land}[3]}
                [\sFmla{\True}{\formula{C}},
                  just={\TRule{\True}{\land}[3]},close
                ]
              ]
            ]
          ]
        ]
        [\sFmla{\False}{\formula{A} \lor \formula{C}},
          just={\TRule{\False}{\land}[2]}, checked
          [\sFmla{\False}{\formula{A}},
            just={\TRule{\False}{\lor}[4]}
            [\sFmla{\False}{\formula{C}},
              just={\TRule{\False}{\lor}[4]}
              [\sFmla{\True}{\formula{B}},
                just={\TRule{\True}{\land}[3]}
                [\sFmla{\True}{\formula{C}},
                  just={\TRule{\True}{\land}[3]},close
                ]
              ]
            ]
          ]
        ]
      ]
    ]
  ]
\end{oltableau}

For comparison, here's a closed !!{tableau} for the same set of
assumptions in which the rules are applied in a different order:
\begin{oltableau}
  [\sFmla{\True}{\formula{A} \lor (\formula{B} \land \formula{C})},
    just=\TAss, checked
    [\sFmla{\False}{(\formula{A} \lor \formula{B}) \land
        (\formula{A} \lor \formula{C})}, just=\TAss, checked
      [\sFmla{\False}{\formula{A} \lor \formula{B}},
        just={\TRule{\False}{\land}[2]}, checked
        [\sFmla{\False}{\formula{A}},
          just={\TRule{\False}{\lor}[3]}
          [\sFmla{\False}{\formula{B}},
            just={\TRule{\False}{\lor}[3]}
            [\sFmla{\True}{\formula{A}},
              just={\TRule{\True}{\lor}[1]},close
            ]
            [\sFmla{\True}{\formula{B} \land \formula{C}},
              just={\TRule{\True}{\lor}[1]}, checked
              [\sFmla{\True}{\formula{B}},
                just={\TRule{\True}{\land}[6]}
                [\sFmla{\True}{\formula{C}},
                  just={\TRule{\True}{\land}[6]},close
                ]
              ]
            ]
          ]
        ]
      ]
      [\sFmla{\False}{\formula{A} \lor \formula{C}},
        just={\TRule{\False}{\land}[2]}, checked
        [\sFmla{\False}{\formula{A}},
          just={\TRule{\False}{\lor}[3]}
          [\sFmla{\False}{\formula{C}},
            just={\TRule{\False}{\lor}[3]}
            [\sFmla{\True}{\formula{A}},
              just={\TRule{\True}{\lor}[1]},close
            ]
            [\sFmla{\True}{\formula{B} \land \formula{C}},
              just={\TRule{\True}{\lor}[1]}, checked
              [\sFmla{\True}{\formula{B}},
                just={\TRule{\True}{\land}[6]}
                [\sFmla{\True}{\formula{C}},
                  just={\TRule{\True}{\land}[6]},close
                ]
              ]
            ]
          ]
        ]
      ]
    ]
  ]
\end{oltableau}
\end{ex}

\begin{prob}
Give closed !!{tableau}s of the following:
\begin{enumerate}
\item $\sFmla{\True}{!A \land (!B \land !C)}, \sFmla{\False}{(!A \land !B) \land !C}$.
\item $\sFmla{\True}{!A \lor (!B \lor !C)}, \sFmla{\False}{(!A \lor !B) \lor !C}$.
\item $\sFmla{\True}{!A \lif (!B \lif !C)}, \sFmla{\False}{!B \lif (!A \lif !C)}$.
\item $\sFmla{\True}{!A}, \sFmla{\False}{\lnot\lnot !A}$.
\end{enumerate}
\end{prob}

\begin{prob}
Give closed !!{tableau}s of the following:
\begin{enumerate}
\item $\sFmla{\True}{(!A \lor !B) \lif !C}, \sFmla{\False}{!A \lif !C}$.
\item $\sFmla{\True}{(!A \lif !C) \land (!B \lif !C)}, \sFmla{\False}{(!A \lor !B) \lif !C}$.
\item $\sFmla{\False}{\lnot(!A \land \lnot !A)}$.
\item $\sFmla{\True}{!B \lif !A}, \sFmla{\False}{\lnot !A \lif \lnot !B}$.
\item $\sFmla{\False}{(!A \lif \lnot !A) \lif \lnot !A}$.
\item $\sFmla{\False}{\lnot(!A \lif !B) \lif \lnot !B}$.
\item $\sFmla{\True}{!A \lif !C}, \sFmla{\False}{\lnot (!A \land \lnot !C)}$.
\item $\sFmla{\True}{!A \land \lnot !C}, \sFmla{\False}{\lnot (!A \lif !C)}$.
\item $\sFmla{\True}{!A \lor !B, \lnot !B}, \sFmla{\False}{!A}$.
\item $\sFmla{\True}{\lnot !A \lor \lnot !B}, \sFmla{\False}{\lnot(!A \land !B)}$.
\item $\sFmla{\False}{(\lnot !A \land \lnot !B) \lif\lnot(!A \lor !B)}$.
\item $\sFmla{\False}{\lnot(!A \lor !B) \lif (\lnot !A \land \lnot !B)}$.
\end{enumerate}
\end{prob}

\begin{prob}
Give closed !!{tableau}s of the following:
\begin{enumerate}
\item $\sFmla{\True}{\lnot(!A \lif !B)}, \sFmla{\False}{!A}$.
\item $\sFmla{\True}{\lnot(!A \land !B)}, \sFmla{\False}{\lnot !A \lor \lnot !B}$.
\item $\sFmla{\True}{!A \lif !B}, \sFmla{\False}{\lnot !A \lor !B}$.
\item $\sFmla{\False}{\lnot \lnot !A \lif !A}$.
\item $\sFmla{\True}{!A \lif !B}, \sFmla{\True}{\lnot !A \lif !B}, \sFmla{\False}{!B}$.
\item $\sFmla{\True}{(!A \land !B) \lif !C}, \sFmla{\False}{(!A \lif !C) \lor (!B \lif !C)}$.
\item $\sFmla{\True}{(!A \lif !B) \lif !A}, \sFmla{\False}{!A}$.
\item $\sFmla{\False}{(!A \lif !B) \lor (!B \lif !C)}$.
\end{enumerate}
\end{prob}




  

\olfileid{fol}{tab}{prq}

\olsection{\usetoken{P}{tableau} with Quantifiers}

\begin{ex}
When dealing with quantifiers, we have to make sure not to violate the
eigenvariable condition, and sometimes this requires us to play around
with the order of carrying out certain inferences. In general, it
helps to try and take care of rules subject to the eigenvariable
condition first (they will be higher up in the finished !!{tableau}).

Let's see how we'd give !!a{tableau} for the !!{sentence}
$\lexists[x][\lnot !A(x)] \lif \lnot \lforall[x][!A(x)]$.
Starting as usual, we start by recording the assumption,
\begin{oltableau}
[\sFmla{\False}{\lexists[x][\lnot \formula{A}(x)]\lif \lnot
    \lforall[x][\formula{A}(x)]}, just=\TAss]
\end{oltableau}
Since the !!{main operator} is $\lif$, we apply the
$\TRule{\False}{\lif}$:
\begin{oltableau}
[\sFmla{\False}{\lexists[x][\lnot \formula{A}(x)]\lif \lnot
    \lforall[x][\formula{A}(x)]}, just=\TAss, checked
  [\sFmla{\True}{\lexists[x][\lnot \formula{A}(x)]},
    just={\TRule{\False}{\lif}[1]}
    [\sFmla{\False}{\lnot\lforall[x][\formula{A}(x)]},
      just={\TRule{\False}{\lif}[1]}]
  ]
]
\end{oltableau}
The next line to deal with is~$2$. We use
$\TRule{\True}{\lexists}$. This requires a new !!{constant}; since no
!!{constant}s yet occur, we can pick any one, say,~$a$.
\begin{oltableau}
[\sFmla{\False}{\lexists[x][\lnot \formula{A}(x)]\lif \lnot
    \lforall[x][\formula{A}(x)]}, just=\TAss, checked
  [\sFmla{\True}{\lexists[x][\lnot \formula{A}(x)]},
    just={\TRule{\False}{\lif}[1]}, checked
    [\sFmla{\False}{\lnot\lforall[x][\formula{A}(x)]},
      just={\TRule{\False}{\lif}[1]}
      [\sFmla{\True}{\lnot\formula{A}(a)}, just={\TRule{\True}{\lexists}[2]}]
    ]
  ]
]
\end{oltableau}
Now we apply $\TRule{\False}{\lnot}$ to line~$3$:
\begin{oltableau}
[\sFmla{\False}{\lexists[x][\lnot \formula{A}(x)]\lif \lnot
    \lforall[x][\formula{A}(x)]}, just=\TAss, checked
  [\sFmla{\True}{\lexists[x][\lnot \formula{A}(x)]},
    just={\TRule{\False}{\lif}[1]}, checked
    [\sFmla{\False}{\lnot\lforall[x][\formula{A}(x)]},
      just={\TRule{\False}{\lif}[1]}, checked
      [\sFmla{\True}{\lnot\formula{A}(a)}, just={\TRule{\True}{\lexists}[2]}
        [\sFmla{\True}{\lforall[x][\formula{A}(x)]},
          just = {\TRule{\False}{\lnot}[3]}
        ]
      ]
    ]
  ]
]
\end{oltableau}
We obtain a closed !!{tableau} by applying $\TRule{\True}{\lnot}$ to
line~$4$, followed by $\TRule{\True}{\lforall}$ to line~$5$.
\begin{oltableau}
[\sFmla{\False}{\lexists[x][\lnot \formula{A}(x)]\lif \lnot
    \lforall[x][\formula{A}(x)]}, just=\TAss, checked
  [\sFmla{\True}{\lexists[x][\lnot \formula{A}(x)]},
    just={\TRule{\False}{\lif}[1]}, checked
    [\sFmla{\False}{\lnot\lforall[x][\formula{A}(x)]},
      just={\TRule{\False}{\lif}[1]}, checked
      [\sFmla{\True}{\lnot\formula{A}(a)}, just={\TRule{\True}{\lexists}[2]}
        [\sFmla{\True}{\lforall[x][\formula{A}(x)]},
          just = {\TRule{\False}{\lnot}[3]}
          [\sFmla{\False}{\formula{A}(a)}, just={\TRule{\True}{\lnot}[4]}
            [\sFmla{\True}{\formula{A}(a)},
              just={\TRule{\True}{\lforall}[5]}, close
            ]
          ]
        ]
      ]
    ]
  ]
]
\end{oltableau}
\end{ex}

\begin{ex}
Let's see how we'd give !!a{tableau} for the set
\[
\sFmla{\False}{\lexists[x][!C(x,b)]},
\sFmla{\True}{\lexists[x][(!A(x) \land !B(x))]},
\sFmla{\True}{\lforall[x][(!B(x) \lif !C(x,b))]}.
\]
Starting as usual, we start with the assumptions:
\begin{oltableau}
  [\sFmla{\False}{\lexists[x][\formula{C}(x,b)]}, just=\TAss
    [\sFmla{\True}{\lexists[x][(\formula{A}(x) \land \formula{B}(x))]},
      just=\TAss
      [\sFmla{\True}{\lforall[x][(\formula{B}(x) \lif \formula{C}(x,b))]},
        just=\TAss]
    ]
  ]
\end{oltableau}
We should always apply a rule with the eigenvariable condition first;
in this case that would be $\TRule{\True}{\lexists}$ to
line~$2$. Since the assumptions contain the !!{constant}~$b$, we have
to use a different one; let's pick~$a$ again.
\begin{oltableau}
  [\sFmla{\False}{\lexists[x][\formula{C}(x,b)]}, just=\TAss
    [\sFmla{\True}{\lexists[x][(\formula{A}(x) \land \formula{B}(x))]},
      just=\TAss, checked
      [\sFmla{\True}{\lforall[x][(\formula{B}(x) \lif \formula{C}(x,b))]},
        just=\TAss
        [\sFmla{\True}{\formula{A}(a) \land \formula{B}(a)},
          just={\TRule{\True}{\lexists}[2]}]
      ]
    ]
  ]
\end{oltableau}
If we now apply $\TRule{\False}{\lexists}$ to line~$1$ or
$\TRule{\True}{\lforall}$ to line~$3$, we have to decide which term~$t$
to substitute for~$x$. Since there is no eigenvariable condition for
these rules, we can pick any term we like. In some cases we may even
have to apply the rule several times with different~$t$s. But as a
general rule, it pays to pick one of the terms already occurring in the
!!{tableau}---in this case, $a$ and~$b$---and in this case we can
guess that $a$ will be more likely to result in a closed branch.
\begin{oltableau}
  [\sFmla{\False}{\lexists[x][\formula{C}(x,b)]}, just=\TAss
    [\sFmla{\True}{\lexists[x][(\formula{A}(x) \land \formula{B}(x))]},
      just=\TAss, checked
      [\sFmla{\True}{\lforall[x][(\formula{B}(x) \lif \formula{C}(x,b))]}, just=\TAss
        [\sFmla{\True}{\formula{A}(a) \land \formula{B}(a)}, just={\TRule{\True}{\lexists}[2]}
          [\sFmla{\False}{\formula{C}(a,b)}, just={\TRule{\False}{\lexists}[1]}
            [\sFmla{\True}{\formula{B}(a) \lif \formula{C}(a,b)},
              just={\TRule{\True}{\lforall}[3]}
            ]
          ]
        ]
      ]
    ]
  ]
\end{oltableau}
We don't check the !!{signed formula}s in lines $1$ and $3$, since we may
have to use them again. Now apply $\TRule{\True}{\land}$ to line~$4$:
\begin{oltableau}
  [\sFmla{\False}{\lexists[x][\formula{C}(x,b)]}, just=\TAss
    [\sFmla{\True}{\lexists[x][(\formula{A}(x) \land \formula{B}(x))]},
      just=\TAss, checked
      [\sFmla{\True}{\lforall[x][(\formula{B}(x) \lif \formula{C}(x,b))]}, just=\TAss
        [\sFmla{\True}{\formula{A}(a) \land \formula{B}(a)},
          just={\TRule{\True}{\lexists}[2]}, checked
          [\sFmla{\False}{\formula{C}(a,b)}, just={\TRule{\False}{\lexists}[1]}
            [\sFmla{\True}{\formula{B}(a) \lif \formula{C}(a,b)},
              just={\TRule{\True}{\lforall}[3]}
              [\sFmla{\True}{\formula{A}(a)}, just={\TRule{\True}{\land}[4]}
                [\sFmla{\True}{\formula{B}(a)}, just={\TRule{\True}{\land}[4]}
                ]
              ]
            ]
          ]
        ]
      ]
    ]
  ]
\end{oltableau}
If we now apply $\TRule{\True}{\lif}$ to line~$6$, the !!{tableau} closes:
\begin{oltableau}
  [\sFmla{\False}{\lexists[x][\formula{C}(x,b)]}, just=\TAss
    [\sFmla{\True}{\lexists[x][(\formula{A}(x) \land \formula{B}(x))]},
      just=\TAss, checked
      [\sFmla{\True}{\lforall[x][(\formula{B}(x) \lif \formula{C}(x,b))]},
        just=\TAss
        [\sFmla{\True}{\formula{A}(a) \land \formula{B}(a)},
          just={\TRule{\True}{\lexists}[2]}, checked
          [\sFmla{\False}{\formula{C}(a,b)},
            just={\TRule{\False}{\lexists}[1]}
            [\sFmla{\True}{\formula{B}(a) \lif \formula{C}(a,b)},
              just={\TRule{\True}{\lforall}[3]}, checked
              [\sFmla{\True}{\formula{A}(a)},
                just={\TRule{\True}{\land}[4]}
                [\sFmla{\True}{\formula{B}(a)},
                  just={\TRule{\True}{\land}[4]}
                  [\sFmla{\False}{\formula{B}(a)},
                    just={\TRule{\True}{\lif}[6]},close
                  ]
                  [\sFmla{\True}{\formula{C}(a,b)},
                    just={\TRule{\True}{\lif}[6]},close
                  ]
                ]
              ]
            ]
          ]
        ]
      ]
    ]
  ]
\end{oltableau}


\end{ex}

\begin{ex}
We construct !!a{tableau} for the set
\[
\sFmla{\True}{\lforall[x][!A(x)]}, \sFmla{\True}{\lforall[x][!A(x)] 
  \lif \lexists[y][!B(y)]}, \sFmla{\True}{\lnot\lexists[y][!B(y)]}.
\]
Starting as usual, we write down the assumptions:
\begin{oltableau}
  [\sFmla{\True}{\lforall[x][\formula{A}(x)]}, just=\TAss
    [\sFmla{\True}{\lforall[x][\formula{A}(x)] \lif
        \lexists[y][\formula{B}(y)]}, just=\TAss
      [\sFmla{\True}{\lnot\lexists[y][\formula{B}(y)]}, just=\TAss
      ]
    ]
  ]
\end{oltableau}
We begin by applying the $\TRule{\True}{\lnot}$ rule to line~$3$. A
corollary to the rule ``always apply rules with eigenvariable
conditions first'' is ``defer applying quantifier rules without
eigenvariable conditions until needed.'' Also, defer rules that result
in a split.
\begin{oltableau}
  [\sFmla{\True}{\lforall[x][\formula{A}(x)]}, just=\TAss
    [\sFmla{\True}{\lforall[x][\formula{A}(x)] \lif
        \lexists[y][\formula{B}(y)]}, just=\TAss
      [\sFmla{\True}{\lnot\lexists[y][\formula{B}(y)]}, just=\TAss, checked
        [\sFmla{\False}{\lexists[y][\formula{B}(y)]}, just={\TRule{\True}{\lnot}[3]}]
      ]
    ]
  ]
\end{oltableau}
The new line~$4$ requires $\TRule{\False}{\lexists}$, a quantifier
rule without the eigenvariable condition. So we defer this in favor of
using $\TRule{\True}{\lif}$ on line~$2$.
\begin{oltableau}
  [\sFmla{\True}{\lforall[x][\formula{A}(x)]}, just=\TAss
    [\sFmla{\True}{\lforall[x][\formula{A}(x)] \lif
        \lexists[y][\formula{B}(y)]}, just=\TAss, checked
      [\sFmla{\True}{\lnot\lexists[y][\formula{B}(y)]}, just=\TAss, checked
        [\sFmla{\False}{\lexists[y][\formula{B}(y)]},
          just={\TRule{\True}{\lnot}[3]},
          [\sFmla{\False}{\lforall[x][\formula{A}(x)]}, just={\TRule{\True}{\lif}[2]}]
          [\sFmla{\True}{\lexists[y][\formula{B}(y)]}, just={\TRule{\True}{\lif}[2]}]
        ]
      ]
    ]
  ]
\end{oltableau}
Both new !!{signed formula}s require rules with eigenvariable conditions, so
these should be next:
\begin{oltableau}
  [\sFmla{\True}{\lforall[x][\formula{A}(x)]}, just=\TAss
    [\sFmla{\True}{\lforall[x][\formula{A}(x)] \lif
        \lexists[y][\formula{B}(y)]}, just=\TAss, checked
      [\sFmla{\True}{\lnot\lexists[y][\formula{B}(y)]}, just=\TAss, checked
        [\sFmla{\False}{\lexists[y][\formula{B}(y)]},
          just={\TRule{\True}{\lnot}[3]}
          [\sFmla{\False}{\lforall[x][\formula{A}(x)]}, just={\TRule{\True}{\lif}[2]},checked
            [\sFmla{\False}{\formula{A}(b)}, just={\TRule{\False}{\lforall}[5]}]
          ]
          [\sFmla{\True}{\lexists[y][\formula{B}(y)]}, just={\TRule{\True}{\lif}[2]},checked
            [\sFmla{\True}{\formula{B}(c)}, just={\TRule{\True}{\lexists}[5]}]
          ]
        ]
      ]
    ]
  ]
\end{oltableau}
To close the branches, we have to use the !!{signed formula}s on lines $1$
and~$3$. The corresponding rules (\TRule{\True}{\lforall} and
\TRule{\False}{\lexists}) don't have eigenvariable conditions, so we
are free to pick whichever terms are suitable. In this case, that's
$b$ and~$c$, respectively.
\begin{oltableau}
  [\sFmla{\True}{\lforall[x][\formula{A}(x)]}, just=\TAss
    [\sFmla{\True}{\lforall[x][\formula{A}(x)] \lif
        \lexists[y][\formula{B}(y)]}, just=\TAss, checked
      [\sFmla{\True}{\lnot\lexists[y][\formula{B}(y)]}, just=\TAss, checked
        [\sFmla{\False}{\lexists[y][\formula{B}(y)]},
          just={\TRule{\True}{\lnot}[3]}
          [\sFmla{\False}{\lforall[x][\formula{A}(x)]},
            just={\TRule{\True}{\lif}[2]},checked
            [\sFmla{\False}{\formula{A}(b)},
              just={\TRule{\False}{\lforall}[5]}
              [\sFmla{\True}{\formula{A}(b)},
                just={\TRule{\True}{\lforall}[1]},close
              ]
            ]
          ]
          [\sFmla{\True}{\lexists[y][\formula{B}(y)]},
            just={\TRule{\True}{\lif}[2]},checked
            [\sFmla{\True}{\formula{B}(c)},
              just={\TRule{\True}{\lexists}[5]}
              [\sFmla{\False}{\formula{B}(c)},
                just={\TRule{\False}{\lexists}[4]},close
              ]
            ]
          ]
        ]
      ]
    ]
  ]
\end{oltableau}
\end{ex}

\begin{prob}
Give closed !!{tableau}s of the following:
\begin{enumerate}
\item $\sFmla{\False}{(\lforall[x][!A(x)] \land \lforall[y][!B(y)])
\lif \lforall[z][(!A(z) \land !B(z))]}$.
\item $\sFmla{\False}{(\lexists[x][!A(x)] \lor \lexists[y][!B(y)])
\lif \lexists[z][(!A(z) \lor !B(z))]}$.
\item $\sFmla{\True}{\lforall[x][(!A(x) \lif !B)]},
\sFmla{\False}{\lexists[y][!A(y)] \lif !B}$.
\item $\sFmla{\True}{\lforall[x][\lnot !A(x)]},
\sFmla{\False}{\lnot\lexists[x][!A(x)]}$.
\item $\sFmla{\False}{\lnot\lexists[x][!A(x)] \lif \lforall[x][\lnot
!A(x)]}$.
\item $\sFmla{\False}{\lnot\lexists[x][\lforall[y][((!A(x,y) \lif
\lnot !A(y,y)) \land (\lnot !A(y,y) \lif !A(x,y)))]]}$.
\end{enumerate}
\end{prob}

\begin{prob}
Give closed !!{tableau}s of the following:
\begin{enumerate}
\item $\sFmla{\False}{\lnot\lforall[x][!A(x)] \lif \lexists[x][\lnot!A(x)]}$.
\item $\sFmla{\True}{(\lforall[x][!A(x)] \lif !B)}, \sFmla{\False}{\lexists[y][(!A(y) \lif !B)]}$.
\item $\sFmla{\False}{\lexists[x][(!A(x) \lif \lforall[y][!A(y)])]}$.
\end{enumerate}
\end{prob}






\olfileid{fol}{tab}{ptn}
      
\olsection{Proof-Theoretic Notions}

\begin{editorial}
This section collects the definitions of the provability relation
and consistency for tableaux.
\end{editorial}

\begin{explain}
Just as we've defined a number of important semantic notions
(validity, entailment, satisfiability), we now define corresponding
\emph{proof-theoretic notions}.  These are not defined by appeal to
satisfaction of !!{sentence}s in !!{structure}s, but by appeal to the
existence of certain closed !!{tableau}x.  It was an important
discovery that these notions coincide.  That they do is the content of
the \emph{soundness} and \emph{completeness theorems}.
\end{explain}


\begin{defn}[Theorems]
A !!{sentence}~$!A$ is a \emph{theorem} if there is a closed
!!{tableau} for~$\sFmla{\False}{!A}$.  We write $\Proves !A$ if $!A$
is a theorem and $\Proves/ !A$ if it is not.
\end{defn}

\begin{defn}[!!^{derivability}]
!!^a{sentence} $!A$ is \emph{!!{derivable} from} a set of
!!{sentence}s~$\Gamma$, $\Gamma \Proves !A$ iff there is a
finite set $\{!B_1, \dots, !B_n\} \subseteq \Gamma$
and a closed !!{tableau} for the set
\[
\{
  \sFmla{\False}{!A},
  \sFmla{\True}{!B_1}, \dots,
  \sFmla{\True}{!B_n}
  \}.
\]
If $!A$ is not !!{derivable} from $\Gamma$ we write $\Gamma \Proves/
!A$.
\end{defn}

\begin{defn}[Consistency]
A set of !!{sentence}s~$\Gamma$ is \emph{inconsistent} iff there is a
finite set $\{!B_1, \dots, !B_n\} \subseteq \Gamma$ and a closed
!!{tableau} for the set
\[
\{
  \sFmla{\True}{!B_1}, \dots,
  \sFmla{\True}{!B_n}  
  \}.
\]
If $\Gamma$ is not inconsistent, we say it is \emph{consistent}.
\end{defn}

\begin{prop}[Reflexivity]
\ollabel{prop:reflexivity}
If $!A \in \Gamma$, then $\Gamma \Proves !A$.
\end{prop}

\begin{proof}
If $!A \in \Gamma$, $\{!A\}$ is a finite subset of~$\Gamma$ and the !!{tableau}
\begin{oltableau}
  [\sFmla{\False}{\formula{A}}, just = \TAss
    [\sFmla{\True}{\formula{A}}, just = \TAss,close]
  ]
\end{oltableau}
is closed.
\end{proof}
  

\begin{prop}[Monotonicity]
\ollabel{prop:monotonicity}
If $\Gamma \subseteq \Delta$ and $\Gamma \Proves !A$, then $\Delta
\Proves !A$.
\end{prop}

\begin{proof}
Any finite subset of~$\Gamma$ is also a finite subset of~$\Delta$.
\end{proof}

\begin{prop}[Transitivity]
\ollabel{prop:transitivity}
If $\Gamma \Proves !A$ and $\{!A\} \cup
\Delta \Proves !B$, then $\Gamma \cup \Delta \Proves !B$.
\end{prop}

\begin{proof}
If $\{!A\} \cup \Delta \Proves !B$, then there is a finite subset $\Delta_0 =
\{!C_1, \dots, !C_n\} \subseteq \Delta$ such that
\begin{align*}
\{\sFmla{\False}{!B}, & \sFmla{\True}{!A}, \sFmla{\True}{!C_1},
\dots, \sFmla{\True}{!C_n}\}
\intertext{has a closed !!{tableau}. If $\Gamma \Proves !A$ then there
  are $!D_1$, \dots, $!D_m \subseteq \Gamma$ such that}
\{\sFmla{\False}{!A}, & \sFmla{\True}{!D_1},
\dots, \sFmla{\True}{!D_m}\}
\end{align*}
has a closed !!{tableau}.

Now consider the !!{tableau} with assumptions
\[
\sFmla{\False}{!B},
\sFmla{\True}{!C_1}, \dots, \sFmla{\True}{!C_n},
\sFmla{\True}{!D_1}, \dots, \sFmla{\True}{!D_m}.
\]
Apply the \Cut{} rule on~$!A$. This generates two branches, one has
$\sFmla{\True}{!A}$ in it, the other $\sFmla{\False}{!A}$. Thus,
on the one branch, all of
\[
\{\sFmla{\False}{!B}, \sFmla{\True}{!A},
\sFmla{\True}{!C_1}, \dots, \sFmla{\True}{!C_n}\}
\]
are available. Since there is a closed !!{tableau} for these
assumptions, we can attach it to that branch; every branch through
$\sFmla{\True}{!A}$ closes. On the other branch, all of
\[
\{\sFmla{\False}{!A}, \sFmla{\True}{!D_1}, \dots,
\sFmla{\True}{!D_m}\}
\]
are available, so we can also complete the other side to obtain a
closed !!{tableau}.  This shows $\Gamma \cup \Delta \Proves !B$.
\end{proof}

Note that this means that in particular if $\Gamma \Proves !A$ and $!A
\Proves !B$, then $\Gamma \Proves !B$. It follows also that if $!A_1,
\dots, !A_n \Proves !B$ and $\Gamma \Proves !A_i$ for each~$i$, then
$\Gamma \Proves !B$.

\begin{prop}
\ollabel{prop:incons}
$\Gamma$ is inconsistent iff $\Gamma \Proves !A$ for every
  !!{sentence}~$!A$.
\end{prop}

\begin{proof}
Exercise.
\end{proof}


\begin{prob}
Prove \olref[fol][tab][ptn]{prop:incons}
\end{prob}




\begin{prop}[Compactness]
  \ollabel{prop:proves-compact}
  \begin{enumerate}
  \item If $\Gamma \Proves !A$ then there is a finite subset $\Gamma_0
    \subseteq \Gamma$ such that $\Gamma_0 \Proves !A$.
  \item If every finite subset of~$\Gamma$ is
    consistent, then $\Gamma$ is consistent.
  \end{enumerate}
\end{prop}

\begin{proof}
  \begin{enumerate}
    \item If $\Gamma \Proves !A$, then there is a finite subset
      $\Gamma_0 = \{!B_1, \dots, !B_n\}$ and a closed !!{tableau} for
      \[
      \{\sFmla{\False}{!A}, \sFmla{\True}{!B_1}, \dots, \sFmla{\True}{!B_n}\}
      \]
      This !!{tableau} also shows $\Gamma_0 \Proves !A$.
    \item If $\Gamma$ is inconsistent, then for some finite subset
      $\Gamma_0 = \{!B_1, \dots, !B_n\}$ there is a closed !!{tableau}
      for
      \[
      \{\sFmla{\True}{!B_1}, \dots, \sFmla{\True}{!B_n}\}
      \]
      This closed !!{tableau} shows that $\Gamma_0$ is inconsistent.
  \end{enumerate}
\end{proof}





\olfileid{fol}{tab}{prv}

\olsection{\usetoken{S}{derivability} and Consistency}

We will now establish a number of properties of the !!{derivability}
relation.  They are independently interesting, but each will play a
role in the proof of the completeness theorem.

\begin{prop}\ollabel{prop:provability-contr}
  If $\Gamma \Proves !A$ and $\Gamma \cup \{!A\}$ is
  inconsistent, then $\Gamma$ is inconsistent.
\end{prop}

\begin{proof}
There are finite $\Gamma_0 = \{!B_1, \dots, !B_n\}$ and $\Gamma_1
  =\{!C_1, \dots, !C_n\} \subseteq \Gamma$ such that
  \begin{align*}
    \{\sFmla{\False}{!A}, &
    \sFmla{\True}{!B_1}, \dots, \sFmla{\True}{!B_n}\} \\
    \{\sFmla{\True}{!A}, &
    \sFmla{\True}{!C_1}, \dots, \sFmla{\True}{!C_m}\}
  \end{align*}
  have closed !!{tableau}s.  Using the \Cut{} rule on $!A$ we can
  combine these into a single closed !!{tableau} that shows $\Gamma_0
  \cup \Gamma_1$ is inconsistent.  Since $\Gamma_0
  \subseteq \Gamma$ and $\Gamma_1 \subseteq \Gamma$, $\Gamma_0 \cup
  \Gamma_1 \subseteq \Gamma$, hence $\Gamma$~is inconsistent.
\end{proof}

\begin{prop}
\ollabel{prop:prov-incons}
$\Gamma \Proves !A$ iff $\Gamma \cup \{\lnot !A\}$ is inconsistent.
\end{prop}

\begin{proof}
First suppose $\Gamma \Proves !A$, i.e., there is
a closed !!{tableau} for
\[
\{\sFmla{\False}{!A},
\sFmla{\True}{!B_1}, \dots, \sFmla{\True}{!B_n}\}
\]
Using the $\TRule{\True}{\lnot}$ rule, this can be turned into a
closed !!{tableau} for
\[
\{\sFmla{\True}{\lnot !A},
\sFmla{\True}{!B_1}, \dots, \sFmla{\True}{!B_n}\}.
\]

On the other hand, if there is a closed !!{tableau} for the latter, we
can turn it into a closed !!{tableau} of the former by removing every
formula that results from \TRule{\True}{\lnot} applied to the first
assumption~$\sFmla{\True}{\lnot !A}$ as well as that assumption, and
adding the assumption $\sFmla{\False}{!A}$. For if a branch was closed
before because it contained the conclusion of \TRule{\True}{\lnot}
applied to $\sFmla{\True}{\lnot !A}$, i.e., $\sFmla{\False}{!A}$, the
corresponding branch in the new !!{tableau} is also closed. If a
branch in the old tableau was closed because it contained the
assumption $\sFmla{\True}{\lnot !A}$ as well as $\sFmla{\False}{\lnot
  !A}$ we can turn it into a closed branch by applying
$\TRule{\False}{\lnot}$ to $\sFmla{\False}{\lnot !A}$ to obtain
$\sFmla{\True}{!A}$. This closes the branch since we added
$\sFmla{\False}{!A}$ as an assumption.
\end{proof}

\begin{prob}
Prove that $\Gamma \Proves \lnot !A$ iff $\Gamma \cup \{!A\}$ is inconsistent.
\end{prob}

\begin{prop}\ollabel{prop:explicit-inc}
  If $\Gamma \Proves !A$ and $\lnot !A \in \Gamma$, then $\Gamma$ is
  inconsistent.
\end{prop}

\begin{proof}
  Suppose $\Gamma \Proves !A$ and $\lnot !A \in \Gamma$.  Then there
  are $!B_1$, \dots, $!B_n \in \Gamma$ such that \
  \[
  \{\sFmla{\False}{!A}, \sFmla{\True}{!B_1}, \dots, \sFmla{\True}{!B_n}\}
  \]
  has a closed tableau. Replace the assumption \sFmla{\False}{!A} by
  \sFmla{\True}{\lnot !A}, and insert the conclusion of
  \TRule{\True}{\lnot} applied to \sFmla{\False}{!A} after the
  assumptions. Any !!{sentence} in the !!{tableau} justified by appeal
  to line~$1$ in the old !!{tableau} is now justified by appeal to
  line~$n+1$. So if the old !!{tableau} was closed, the new one is.
  It shows that $\Gamma$ is inconsistent, since all assumptions are
  in~$\Gamma$.
\end{proof}

\begin{prop}\ollabel{prop:provability-exhaustive}
  If $\Gamma \cup \{!A\}$ and $\Gamma \cup \{\lnot !A\}$ are both
  inconsistent, then $\Gamma$ is inconsistent.
\end{prop}

\begin{proof}
  If there are $!B_1$, \dots, $!B_n \in \Gamma$ and $!C_1$, \dots,
  $!C_m \in \Gamma$ such that
  \begin{align*}
    \{\sFmla{\True}{!A}, &
    \sFmla{\True}{!B_1}, \dots, \sFmla{\True}{!B_n}\} \text{ and}\\
    \{\sFmla{\True}{\lnot !A}, &
    \sFmla{\True}{!C_1}, \dots, \sFmla{\True}{!C_m}\}
  \end{align*}
  both have closed !!{tableau}s, we can construct a single, combined
  !!{tableau} that shows that $\Gamma$ is inconsistent by using as
  assumptions $\sFmla{\True}{!B_1}$, \dots, $\sFmla{\True}{!B_n}$
  together with $\sFmla{\True}{!C_1}$, \dots, $\sFmla{\True}{!C_m}$,
  followed by an application of the \Cut{} rule. This yields two
  branches, one starting with $\sFmla{\True}{!A}$, the other with
  $\sFmla{\False}{!A}$.  
  
  On the left left side, add the part of the first
  !!{tableau} below its assumptions. Here, every rule application is still
  correct, since each of the assumptions of the first !!{tableau},
  including $\sFmla{\True}{!A}$, is available. Thus, every branch
  below $\sFmla{\True}{!A}$ closes. 
  
  On the right side, add the part of the second !!{tableau} below its
  assumption, with the results of any applications
  of~$\TRule{\True}{\lnot}$ to $\sFmla{\True}{\lnot !A}$ removed.  The
  conclusion of $\TRule{\True}{\lnot}$ to $\sFmla{\True}{\lnot !A}$ is
  $\sFmla{\False}{!A}$, which is nevertheless available, as it is the
  conclusion of the \Cut{} rule on the right side of the combined !!{tableau}.

  If a branch in the second tableau was closed because it contained
  the assumption $\sFmla{\True}{\lnot !A}$ (which no longer appears as
  an assumption in the combined !!{tableau}) as well as
  $\sFmla{\False}{\lnot !A}$, we can applying $\TRule{\False}{\lnot}$
  to $\sFmla{\False}{\lnot !A}$ to obtain $\sFmla{\True}{!A}$. Now the
  corresponding branch in the combined !!{tableau} also closes,
  because it contains the right-hand conclusion of the \Cut{} rule,
  $\sFmla{\False}{!A}$. If a branch in the second !!{tableau} closed
  for any other reason, the corresponding branch in the combined
  !!{tableau} also closes, since any !!{signed formula}s other than
  $\sFmla{\True}{\lnot !A}$ occurring on the branch in the old, second
  !!{tableau} also occur on the corresponding branch in the combined
  !!{tableau}.
  \end{proof}





\olfileid{fol}{tab}{ppr}

\olsection{\usetoken{S}{derivability} and the Propositional Connectives}

\begin{explain}
  We establish that the !!{derivability} relation~$\Proves$ of
  tableaux is strong enough to establish some basic facts involving
  the propositional connectives, such as that $!A \land !B \Proves !A$
  and $!A, !A \lif !B \Proves !B$ (modus ponens). These facts are
  needed for the proof of the completeness theorem.
\end{explain}

\begin{prop}\ollabel{prop:provability-land}
  \begin{enumerate}
  \item \ollabel{prop:provability-land-left} Both $!A \land !B \Proves
    !A$ and $!A \land !B \Proves !B$.
  \item \ollabel{prop:provability-land-right} $!A, !B \Proves !A \land
    !B$.
  \end{enumerate}
\end{prop}

\begin{proof}
  \begin{enumerate}
  \item Both $\{\sFmla{\False}{!A}, \sFmla{\True}{!A \land !B}\}$ and
    $\{\sFmla{\False}{!B}, \sFmla{\True}{!A \land !B}\}$ have closed
    !!{tableau}s
    \begin{oltableau}{}
      [\sFmla{\False}{\formula{A}}, just=\TAss
        [\sFmla{\True}{\formula{A} \land \formula{B}}, just=\TAss
          [\sFmla{\True{\formula{A}}},just={\TRule{\True}{\land}[2]}
            [\sFmla{\True{\formula{B}}},just={\TRule{\True}{\land}[2]}, close
            ]
          ]
        ]
      ]
    \end{oltableau}
    \begin{oltableau}{}
      [\sFmla{\False}{\formula{B}}, just=\TAss
        [\sFmla{\True}{\formula{A} \land \formula{B}}, just=\TAss
          [\sFmla{\True{\formula{A}}},just={\TRule{\True}{\land}[2]}
            [\sFmla{\True{\formula{B}}},just={\TRule{\True}{\land}[2]}, close
            ]
          ]
        ]
      ]
    \end{oltableau}
    \item Here is a closed !!{tableau} for $\{\sFmla{\True}{!A},
      \sFmla{\True}{!B}, \sFmla{\False}{!A \land !B}\}$:
      \begin{oltableau}
        [\sFmla{\False}{\formula{A} \land \formula{B}}, just = \TAss
          [\sFmla{\True}{\formula{A}}, just = \TAss
            [\sFmla{\True}{\formula{B}}, just=\TAss
              [\sFmla{\False}{\formula{A}}, just = {\TRule{\False}{\land}[1]}, close]
              [\sFmla{\False}{\formula{B}}, just = {\TRule{\False}{\land}[1]}, close]
            ]
          ]
        ]
      \end{oltableau}     
  \end{enumerate}
\end{proof}

\begin{prop}\ollabel{prop:provability-lor}
  \begin{enumerate}
  \item $\{!A \lor !B, \lnot !A, \lnot !B\}$ is inconsistent.
  \item Both $!A \Proves !A \lor !B$ and $!B \Proves !A \lor !B$.
  \end{enumerate}
\end{prop}

\begin{proof}
  \begin{enumerate}
  \item We give a closed !!{tableau} of $\{\sFmla{\True}{!A \lor !B},
    \sFmla{\True}{\lnot !A}, \sFmla{\True}{\lnot !B}\}$:
    \begin{oltableau}
      [\sFmla{\True}{\formula{A} \lor \formula{B}}, just = \TAss
        [\sFmla{\True}{\lnot \formula{A}}, just = \TAss
          [\sFmla{\True}{\lnot \formula{B}}, just = \TAss
            [\sFmla{\False}{\formula{A}}, just = {\TRule{\True}{\lnot}[2]}
              [\sFmla{\False}{\formula{B}}, just = {\TRule{\True}{\lnot}[3]}
                [\sFmla{\True}{\formula{A}}, just = {\TRule{\True}{\lor}[1]}, close]
                [\sFmla{\True}{\formula{B}}, just = {\TRule{\True}{\lor}[1]}, close]
              ]
            ]
          ]
        ]
      ]
    \end{oltableau}
  \item Both $\{\sFmla{\False}{!A \lor !B}, \sFmla{\True}{!A}\}$ and
    $\{\sFmla{\False}{!A \lor !B}, \sFmla{\True}{!B}\}$ have closed
    !!{tableau}s:
    \begin{oltableau}{}
      [\sFmla{\False}{\formula{A} \lor \formula{B}}, just=\TAss
        [\sFmla{\True}{\formula{A}}, just=\TAss
          [\sFmla{\False{\formula{A}}},just={\TRule{\False}{\lor}[1]}
            [\sFmla{\False{\formula{B}}},just={\TRule{\False}{\lor}[1]}, close
            ]
          ]
        ]
      ]
    \end{oltableau}
    \begin{oltableau}{}
      [\sFmla{\False}{\formula{A} \lor \formula{B}}, just=\TAss
        [\sFmla{\True}{\formula{B}}, just=\TAss
          [\sFmla{\False{\formula{A}}},just={\TRule{\False}{\lor}[1]}
            [\sFmla{\False{\formula{B}}},just={\TRule{\False}{\lor}[1]}, close
            ]
          ]
        ]
      ]
    \end{oltableau}
  \end{enumerate}
\end{proof}

\begin{prop}\ollabel{prop:provability-lif}
  \begin{enumerate}
  \item \ollabel{prop:provability-lif-left}  $!A, !A \lif !B \Proves !B$.
  \item \ollabel{prop:provability-lif-right}
    Both $\lnot !A \Proves !A \lif !B$ and $!B \Proves !A \lif !B$.
  \end{enumerate}
\end{prop}

\begin{proof}
  \begin{enumerate}
    \item $\{\sFmla{\False}{!B}, \sFmla{\True}{!A \lif !B},
      \sFmla{\True}{!A}\}$ has a closed !!{tableau}:
      \begin{oltableau}
        [\sFmla{\False}{\formula{B}}, just=\TAss
          [\sFmla{\True}{\formula{A} \lif \formula{B}}, just=\TAss
            [\sFmla{\True}{\formula{A}}, just=\TAss
              [\sFmla{\False}{\formula{A}}, just = {\TRule{\True}{\lif}[2]}, close]
              [\sFmla{\True}{\formula{B}}, just = {\TRule{\True}{\lif}[2]}, close]
            ]
          ]
        ]
      \end{oltableau}
    \item Both $\{\sFmla{\False}{!A \lif !B},
      \sFmla{\True}{\lnot !A}\}$ and $\{\sFmla{\False}{!A \lif !B},
      \sFmla{\True}{!B}\}$ have closed !!{tableau}s:
      \begin{oltableau}
        [\sFmla{\False}{\formula{A} \lif \formula{B}}, just = \TAss
          [\sFmla{\True}{\lnot \formula{A}}, just = \TAss
            [\sFmla{\True}{\formula{A}}, just = {\TRule{\False}{\lif}[1]}
              [\sFmla{\False}{\formula{B}}, just = {\TRule{\False}{\lif}[1]}
                [\sFmla{\False}{\formula{A}}, just = {\TRule{\True}{\lnot}[2]}, close]
              ]
            ]
          ]
        ]
      \end{oltableau}
      \begin{oltableau}
        [\sFmla{\False}{\formula{A} \lif \formula{B}}, just = \TAss
          [\sFmla{\True}{\formula{B}}, just = \TAss
            [\sFmla{\True}{\formula{A}}, just = {\TRule{\False}{\lif}[1]}
              [\sFmla{\False}{\formula{B}}, just = {\TRule{\False}{\lif}[1]},close
              ]
            ]
          ]
        ]
      \end{oltableau}
  \end{enumerate}
\end{proof}




  

\olfileid{fol}{tab}{qpr}
\olsection{\usetoken{S}{derivability} and the Quantifiers}

\begin{explain}
  The completeness theorem also requires that the tableaux rules yield
  the facts about~$\Proves$ established in this section.
\end{explain}

\begin{thm}
\ollabel{thm:strong-generalization} If $c$ is a constant not occurring
in $\Gamma$ or~$!A(x)$ and $\Gamma \Proves !A(c)$, then $\Gamma
\Proves \lforall[x][!A(x)]$.
\end{thm}

\begin{proof}
Suppose $\Gamma \Proves !A(c)$, i.e., there are $!B_1$, \dots, $!B_n
\in \Gamma$ and a closed !!{tableau} for
\begin{align*}
\{ \sFmla{\False}{!A(c)}, &
\sFmla{\True}{!B_1}, \dots, \sFmla{\True}{!B_n} \}.
\intertext{We have to show that there is also a closed !!{tableau} for}
\{ \sFmla{\False}{\lforall[x][!A(x)]}, &
\sFmla{\True}{!B_1}, \dots, \sFmla{\True}{!B_n} \}.
\end{align*}
Take the closed !!{tableau} and replace the first assumption with
$\sFmla{\False}{\lforall[x][!A(x)]}$, and insert
$\sFmla{\False}{!A(c)}$ after the assumptions.
\begin{center}
\begin{tableau}{not line numbering}
  [\sFmla{\False}{\formula{A}(c)}
    [\sFmla{\True}{\formula{B}_1}
      [\vdots
      [\sFmla{\True}{\formula{B}_n}, 
        [][][]]]]]
\end{tableau}\qquad
\begin{tableau}{not line numbering}
  [\sFmla{\False}{\lforall[x][\formula{A}(x)]}
    [\sFmla{\True}{\formula{B}_1}
      [\vdots
        [\sFmla{\True}{\formula{B}_n},
          [\sFmla{\False}{\formula{A}(c)}
        [][][]]]]]]
\end{tableau}
\end{center}
The tableau is still closed, since all !!{sentence}s available as
assumptions before are still available at the top of the
!!{tableau}. The inserted line is the result of a correct application
of $\TRule{\False}{\lforall}$, since the !!{constant}~$c$ does not
occur in $!B_1$, \dots, $!B_n$ or $\lforall[x][!A(x)]$, i.e., it does
not occur above the inserted line in the new !!{tableau}.
\end{proof}

\begin{prop}
\ollabel{prop:provability-quantifiers}
\begin{tagenumerate}{prvEx,prvAll}
$!A(t) \Proves \lexists[x][!A(x)]$.
$\lforall[x][!A(x)] \Proves !A(t)$.
\end{tagenumerate}
\end{prop}

\begin{proof}
\begin{tagenumerate}{prvEx,prvAll}
A closed !!{tableau} for
  $\sFmla{\False}{\lexists[x][!A(x)]}, \sFmla{\True}{!A(t)}$ is:
  \begin{oltableau}
    [\sFmla{\False}{\lexists[x][\formula{A}(x)]}, just = \TAss
      [\sFmla{\True}{\formula{A}(t)}, just = \TAss
        [\sFmla{\False}{\formula{A}(t)}, just = {\TRule{\False}{\lexists}[1]}, close]]]
    \end{oltableau}
A closed !!{tableau} for
  $\sFmla{\False}{\formula{A}(t)}, \sFmla{\True}{\lforall[x][!A(x)]}, $ is:
  \begin{oltableau}
    [\sFmla{\False}{\formula{A}(t)}, just = \TAss
      [\sFmla{\True}{\lforall[x][\formula{A}(x)]}, just = \TAss
        [\sFmla{\True}{\formula{A}(t)}, just = {\TRule{\True}{\lforall}[2]}, close]]]
    \end{oltableau}
\end{tagenumerate}
\end{proof}






\olfileid{fol}{tab}{sou}
      
\olsection{Soundness}

\begin{explain}
!!^a{derivation} system, such as tableaux, is \emph{sound}
if it cannot !!{derive} things that do not actually hold.  Soundness is
thus a kind of guaranteed safety property for !!{derivation} systems.
Depending on which proof theoretic property is in question, we would
like to know for instance, that
\begin{enumerate}
\item every !!{derivable}~$!A$ is valid;
\item if !!a{sentence} is !!{derivable} from some others, it is also a
  consequence of them;
\item if a set of !!{sentence}s is inconsistent, it is unsatisfiable.
\end{enumerate}
These are important properties of !!a{derivation} system.  If any of them do
not hold, the !!{derivation} system is deficient---it would !!{derive} too much.
Consequently, establishing the soundness of !!a{derivation} system is of the
utmost importance.

Because all these proof-theoretic properties are defined via closed
!!{tableau}s of some kind or other, proving (1)--(3) above requires
proving something about the semantic properties of closed
!!{tableau}s.  We will first define what it means for !!a{signed
  formula} to be satisfied in a structure, and then show that if a
!!{tableau} is closed, no structure satisfies all its assumptions.
(1)--(3) then follow as corollaries from this result.
\end{explain}

\begin{defn}
!!^a{structure}~$\Struct{M}$
\emph{satisfies} !!a{signed formula} $\sFmla{\True}{!A}$ iff
$\Sat{M}{!A}$, and it satisfies
$\sFmla{\False}{!A}$ iff
$\Sat/{M}{!A}$. $\Struct{M}$
satisfies a set of !!{signed formula}s~$\Gamma$ iff it satisfies every
$\sFmla{S}{!A} \in \Gamma$. $\Gamma$~is \emph{satisfiable} if there is
!!a{structure} that satisfies it, and
\emph{unsatisfiable} otherwise.
\end{defn}

\begin{thm}[Soundness]
  \ollabel{thm:tableau-soundness}
  If $\Gamma$ has a closed !!{tableau}, $\Gamma$ is unsatisfiable.
\end{thm}

\begin{proof}
Let's call a branch of !!a{tableau} satisfiable iff the set of
!!{signed formula}s on it is satisfiable, and let's call !!a{tableau}
satisfiable if it contains at least one satisfiable branch.

We show the following: Extending a satisfiable !!{tableau} by one of
the rules of inference always results in a satisfiable !!{tableau}.
This will prove the theorem: any closed !!{tableau} results by
applying rules of inference to the !!{tableau} consisting only of
assumptions from~$\Gamma$. So if $\Gamma$ were satisfiable, any
!!{tableau} for it would be satisfiable. A closed !!{tableau},
however, is clearly not satisfiable: every branch contains both
$\sFmla{\True}{!A}$ and $\sFmla{\False}{!A}$, and no structure can
both satisfy and not satisfy~$!A$.

Suppose we have a satisfiable !!{tableau}, i.e., !!a{tableau} with at
least one satisfiable branch. Applying a rule of inference either adds
!!{signed formula}s to a branch, or splits a branch in two. If the
!!{tableau} has a satisfiable branch which is not extended by the rule
application in question, it remains a satisfiable branch in the
extended !!{tableau}, so the extended tableau is satisfiable. So we
only have to consider the case where a rule is applied to a
satisfiable branch.

Let $\Gamma$ be the set of !!{signed formula}s on that branch, and
let $\sFmla{S}{!A} \in \Gamma$ be the !!{signed formula} to which the
rule is applied. If the rule does not result in a split branch, we
have to show that the extended branch, i.e., $\Gamma$ together with
the conclusions of the rule, is still satisfiable. If the rule results
in a split branch, we have to show that at least one of the two
resulting branches is satisfiable.
  
First, we consider the possible inferences that do not result in a split branch.
\begin{enumerate}
\item The branch is expanded by applying $\TRule{\True}{\lnot}$ to
  $\sFmla{\True}{\lnot !B} \in \Gamma$. Then the extended branch
  contains the !!{signed formula}s $\Gamma \cup
  \{\sFmla{\False}{!B}\}$.  Suppose
  $\Sat{M}{\Gamma}$. In particular,
  $\Sat{M}{\lnot !B}$. Thus,
  $\Sat/{M}{!B}$, i.e.,
  $\Struct{M}$ satisfies
  $\sFmla{\False}{!B}$.
\item The branch is expanded by applying $\TRule{\False}{\lnot}$ to
  $\sFmla{\False}{\lnot !B} \in \Gamma$: Exercise.
\item The branch is expanded by applying $\TRule{\True}{\land}$ to
  $\sFmla{\True}{!B \land !C} \in \Gamma$, which results in two new
  !!{signed formula}s on the branch: $\sFmla{\True}{!B}$ and
  $\sFmla{\True}{!C}$. Suppose
  $\Sat{M}{\Gamma}$, in particular
  $\Sat{M}{!B \land !C}$. Then
  $\Sat{M}{!B}$ and
  $\Sat{M}{!C}$. This means that
  $\Struct{M}$ satisfies both
  $\sFmla{\True}{!B}$ and $\sFmla{\True}{!C}$.
\item The branch is expanded by applying $\TRule{\False}{\lor}$ to
  $\sFmla{\False}{!B \lor !C} \in \Gamma$: Exercise.
\item The branch is expanded by applying $\TRule{\False}{\lif}$ to
  $\sFmla{\False}{!B \lif !C} \in \Gamma$: This results in two new
  !!{signed formula}s on the branch: $\sFmla{\True}{!B}$ and
  $\sFmla{\False}{!C}$. Suppose
  $\Sat{M}{\Gamma}$, in particular
  $\Sat/{M}{!B \lif !C}$. Then
  $\Sat{M}{!B}$ and
  $\Sat/{M}{!C}$. This means that
  $\Struct{M}$ satisfies both
  $\sFmla{\True}{!B}$ and $\sFmla{\False}{!C}$.


\item The branch is expanded by applying $\TRule{\True}{\lforall}$ to
  $\sFmla{\True}{\lforall[x][!B(x)]} \in \Gamma$: This results in a
  new !!{signed formula}~$\sFmla{\True}{!A(t)}$ on the branch.
  Suppose $\Sat{M}{\Gamma}$, in particular,
  $\Sat{M}{\lforall[x][!A(x)]}$.  By
  \olref[syn][sem]{prop:quant-terms}, $\Sat{M}{!A(t)}$.  Consequently,
  $\Struct{M}$ satisfies $\sFmla{\True}{!A(t)}$.

\item The branch is expanded by applying $\TRule{\False}{\lforall}$ to
  $\sFmla{\False}{\lforall[x][!B(x)]} \in \Gamma$: This results in a
  new !!{signed formula}~$\sFmla{\False}{!A(a)}$ where $a$ is
  !!a{constant} not occurring in~$\Gamma$. Since $\Gamma$ is
  satisfiable, there is a $\Struct{M}$ such that $\Sat{M}{\Gamma}$, in
  particular $\Sat/{M}{\lforall[x][!B(x)]}$. We have to show that
  $\Gamma \cup \{\sFmla{\False}{!A(a)}\}$ is satisfiable. To do this,
  we define a suitable~$\Struct{M'}$ as follows.

  By \olref[syn][ass]{prop:sat-quant}, $\Sat/{M}{\lforall[x][!B(x)]}$
  iff for some $s$, $\Sat/{M}{!B(x)}[s]$. Now let $\Struct{M'}$ be
  just like $\Struct{M}$, except $\Assign{a}{M'} = s(x)$.  By
  \olref[syn][ext]{cor:extensionality-sent}, for any
  $\sFmla{\True}{!C} \in \Gamma$, $\Sat{M'}{!C}$, and for any
  $\sFmla{\False}{!C} \in \Gamma$, $\Sat/{M'}{!C}$, since $a$ does not
  occur in~$\Gamma$.

  By \olref[syn][ext]{prop:extensionality}, $\Sat/{M'}{!A(x)}[s]$.  By
  \olref[syn][ext]{prop:ext-formulas}, $\Sat/{M'}{!A(a)}[s]$.  Since
  $!A(a)$ is !!a{sentence}, by
  \olref[syn][ass]{prop:sentence-sat-true}, $\Sat/{M'}{!A(a)}$, i.e.,
  $\Struct{M'}$ satisfies $\sFmla{\False}{!A(a)}$.
\item The branch is expanded by applying $\TRule{\True}{\lexists}$ to
  $\sFmla{\True}{\lexists[x][!B(x)]} \in \Gamma$: Exercise.
\item The branch is expanded by applying $\TRule{\False}{\lexists}$ to
  $\sFmla{\False}{\lexists[x][!B(x)]} \in \Gamma$: Exercise.
\end{enumerate}
Now let's consider the possible inferences that result in a split branch.
\begin{enumerate}
\item The branch is expanded by applying $\TRule{\False}{\land}$ to
  $\sFmla{\False}{!B \land !C} \in \Gamma$, which results in two
  branches, a left one continuing through $\sFmla{\False}{!B}$ and a
  right one through $\sFmla{\False}{!C}$. Suppose
  $\Sat{M}{\Gamma}$, in particular
  $\Sat/{M}{!B \land !C}$.  Then
  $\Sat/{M}{!B}$ or
  $\Sat/{M}{!C}$. In the former case,
  $\Struct{M}$ satisfies
  $\sFmla{\False}{!B}$, i.e., $\Struct{M}$
  satisfies the formulas on the left branch. In the latter,
  $\Struct{M}$ satisfies
  $\sFmla{\False}{!C}$, i.e., $\Struct{M}$
  satisfies the formulas on the right branch.
\item The branch is expanded by applying $\TRule{\True}{\lor}$ to
  $\sFmla{\True}{!B \lor !C} \in \Gamma$: Exercise.
\item The branch is expanded by applying $\TRule{\True}{\lif}$ to
  $\sFmla{\True}{!B \lif !C} \in \Gamma$: Exercise.
\item The branch is expanded by \Cut: This results in two branches,
  one containing $\sFmla{\True}{!B}$, the other containing
  $\sFmla{\False}{!B}$. Since $\Sat{M}{\Gamma}$
  and either $\Sat{M}{!B}$ or
  $\Sat/{M}{!B}$,
  $\Struct{M}$ satisfies either the left or
  the right branch.
\end{enumerate}
\end{proof}


\begin{prob}
Complete the proof of \olref[fol][tab][sou]{thm:tableau-soundness}.
\end{prob}




\begin{cor}
\ollabel{cor:weak-soundness}
If $\Proves !A$ then $!A$ is valid.
\end{cor}

\begin{cor}
\ollabel{cor:entailment-soundness}
If $\Gamma \Proves !A$ then $\Gamma \Entails !A$.
\end{cor}

\begin{proof}
  If $\Gamma \Proves !A$ then for some $!B_1$, \dots, $!B_n \in
  \Gamma$, $\{\sFmla{\False}{!A}, \sFmla{\True}{!B_1}, \dots,
  \sFmla{\True}{!B_n}\}$ has a closed !!{tableau}. By
  \olref{thm:tableau-soundness}, every
  !!{structure}~$\Struct{M}$
  either makes some $!B_i$ false or makes $!A$ true.  Hence, if
  $\Sat{M}{\Gamma}$ then also
  $\Sat{M}{!A}$.
\end{proof}

\begin{cor}
\ollabel{cor:consistency-soundness}
If $\Gamma$ is satisfiable, then it is consistent.
\end{cor}

\begin{proof}
We prove the contrapositive.  Suppose that $\Gamma$ is not consistent.
Then there are $!B_1$, \dots, $!B_n \in \Gamma$ and a closed
!!{tableau} for $\{\sFmla{\True}{!B_1}, \dots, \sFmla{\True}{!B_n}\}$.  By
\olref{thm:tableau-soundness}, there is no
$\Struct{M}$ such that
$\Sat{M}{!B_i}$ for all $i=1$, \dots,~$n$.  But
then $\Gamma$ is not satisfiable.
\end{proof}




  

\olfileid{fol}{tab}{ide}

\olsection{\usetoken{P}{tableau} with \usetoken{S}{identity}}

!!^{tableau}s with !!{identity} require additional inference rules.
The rules for $\eq$ are ($t$, $t_1$, and $t_2$ are closed terms):

\begin{defish}
\AxiomC{}
\RightLabel{$\eq$}
\UnaryInfC{\sFmla{\True}{\eq[t][t]}}
\DisplayProof
\hfill
\AxiomC{\sFmla{\True}{\eq[t_1][t_2]}}
\noLine
\UnaryInfC{\sFmla{\True}{!A(t_1)}}
\RightLabel{$\TRule{\True}{\eq}$}
\UnaryInfC{\sFmla{\True}{!A(t_2)}}
\DisplayProof
\hfill
\AxiomC{\sFmla{\True}{\eq[t_1][t_2]}}
\noLine
\UnaryInfC{\sFmla{\False}{!A(t_1)}}
\RightLabel{$\TRule{\False}{\eq}$}
\UnaryInfC{\sFmla{\False}{!A(t_2)}}
\DisplayProof
\end{defish}
Note that in contrast to all the other rules, $\TRule{\True}{\eq}$ and
$\TRule{\False}{\eq}$ require that \emph{two} signed !!{formula}s
already appear on the branch, namely both $\sFmla{\True}{\eq[t_1][t_2]}$
and $\sFmla{S}{!A(t_1)}$.

\begin{ex}
If $s$ and $t$ are closed terms, then $\eq[s][t], !A(s)
\Proves !A(t)$:
\begin{oltableau}
  [\sFmla{\False}{\formula{A}(t)}, just = \TAss
    [\sFmla{\True}{\eq[s][t]}, just = \TAss
      [\sFmla{\True}{\formula{A}(s)}, just = \TAss
        [\sFmla{\True}{\formula{A}(t)}, just={\TRule{\True}{\eq}[2, 3]}, close]
      ]
    ]
  ]
\end{oltableau}
This may be familiar as the principle of substitutability of
identicals, or Leibniz' Law.

!!^{tableau}s prove that $\eq$ is symmetric, i.e., that $\eq[s_1][s_2]
\Proves \eq[s_2][s_1]$:
\begin{oltableau}
  [\sFmla{\False}{\eq[s_2][s_1]}, just = \TAss
    [\sFmla{\True}{\eq[s_1][s_2]}, just = \TAss
      [\sFmla{\True}{\eq[s_1][s_1]}, just = {$\eq$}
        [\sFmla{\True}{\eq[s_2][s_1]}, just = {\TRule{\True}{\eq}[2, 3]}, close]
      ]
    ]
  ]
\end{oltableau}
Here, line $2$ is the first prerequisite
!!{formula}~$\sFmla{\True}{\eq[s_1][s_2]}$ of $\TRule{\True}{\eq}$.
Line~$3$ is the second one, of the form $\sFmla{\True}{!A(s_2)}$---think of
$!A(x)$ as
$\eq[x][s_1]$, then $!A(s_1)$ is $\eq[s_1][s_1]$ and $!A(s_2)$ is $\eq[s_2][s_1]$.

They also prove that $\eq$ is transitive, i.e., that $\eq[s_1][s_2],
\eq[s_2][s_3] \Proves \eq[s_1][s_3]$:
\begin{oltableau}
  [\sFmla{\False}{\eq[s_1][s_3]}, just = \TAss
    [\sFmla{\True}{\eq[s_1][s_2]}, just = \TAss
      [\sFmla{\True}{\eq[s_2][s_3]}, just = \TAss
        [\sFmla{\True}{\eq[s_1][s_3]}, just = {\TRule{\True}{\eq}[3, 2]}, close]
      ]
    ]
  ]
\end{oltableau}
In this !!{tableau}, the first prerequisite !!{formula} of
$\TRule{\True}{\eq}$ is line~$3$, $\sFmla{\True}{\eq[s_2][s_3]}$
($s_2$ plays the role of~$t_1$, and $s_3$ the role of~$t_2$). The
second prerequisite, of the form~$\sFmla{\True}{!A(s_2)}$ is line~$2$.
Here, think of $!A(x)$ as $\eq[s_1][x]$; that makes $!A(s_2)$ into
$\eq[t_1][t_2]$ (i.e., line~$2$) and $!A(s_3)$ into the
!!{formula}~$\eq[s_1][s_3]$ in the conclusion.
\end{ex}

\begin{prob}
Give closed !!{tableau}s for the following:
\begin{enumerate}
\item $\sFmla{\False}{\lforall[x][\lforall[y][((x = y \land !A(x))
      \lif !A(y))]]}$
\item $\sFmla{\False}{\lexists[x][(!A(x) \land
    \lforall[y][(!A(y) \lif y = x)])]}$,\\
  $\sFmla{\True}{\lexists[x][!A(x)] \land
    \lforall[y][\lforall[z][((!A(y) \land !A(z)) \lif y = z)]]}$
\end{enumerate}
\end{prob}


  

\olfileid{fol}{tab}{sid}

\olsection{Soundness with \usetoken{S}{identity}}

\begin{prop}
!!^{tableau}s with rules for identity are sound: no closed !!{tableau}
  is satisfiable.
\end{prop}

\begin{proof}
We just have to show as before that if !!a{tableau} has a satisfiable
branch, the branch resulting from applying one of the rules for~$\eq$
to it is also satisfiable. Let $\Gamma$ be the set of signed
!!{formula}s on the branch, and let $\Struct{M}$ be !!a{structure}
satisfying~$\Gamma$.

Suppose the branch is expanded using~$\eq$, i.e., by adding the signed
!!{formula}~$\sFmla{\True}{\eq[t][t]}$. Trivially,
$\Sat{M}{\eq[t][t]}$, so $\Struct{M}$ also satisfies $\Gamma \cup
\{\sFmla{\True}{\eq[t][t]}\}$.

If the branch is expanded using $\TRule{\True}{\eq}$, we add a signed
!!{formula}~$\sFmla{S}{!A(t_2)}$, but $\Gamma$ contains
both~$\sFmla{\True}{\eq[t_1][t_2]}$ and $\sFmla{\True}{!A(t_1)}$. Thus
we have $\Sat{M}{\eq[t_1][t_2]}$ and $\Sat{M}{!A(t_1)}$. Let $s$ be a
variable assignment with $s(x) = \Value{t_1}{M}$.  By
\olref[syn][ass]{prop:sentence-sat-true}, $\Sat{M}{!A(t_1)}[s]$. Since
$\varAssign{s}{s}{x}$, by \olref[syn][ext]{prop:ext-formulas},
$\Sat{M}{!A(x)}[s]$. since $\Sat{M}{\eq[t_1][t_2]}$, we have
$\Value{t_1}{M} = \Value{t_2}{M}$, and hence $s(x) = \Value{t_2}{M}$.
By applying \olref[syn][ext]{prop:ext-formulas} again, we also have
$\Sat{M}{!A(t_2)}[s]$.  By \olref[syn][ass]{prop:sentence-sat-true},
$\Sat{M}{!A(t_2)}$. The case of $\TRule{\False}{\eq}$ is treated
similarly.
\end{proof}




\OLEndChapterHook





  

